\chapter[Wavelets and spectral inference]{Time--frequency representations and estimator uncertainty}
\label{ch:spectral-inference}
A discrete Fourier transform can be computed exactly while
an estimate of an unknown stochastic power spectrum remains
noisy. The distinction is central to electrophysiological inference.

\section{Why the raw periodogram is not generally consistent}
For a zero-mean Gaussian white-noise series with variance $\sigma^2$
and length $T$, define at a non-DC Fourier frequency
\[
I_T(f_k)=\frac{1}{T}\left|\sum_{t=0}^{T-1}
 X_t e^{-2\pi ikt/T}\right|^2.
\]
At generic positive Fourier frequencies below Nyquist, real and
imaginary Fourier coefficients are independent normals with the
appropriate variance; hence
$2I_T(f_k)/\sigma^2\sim\chi^2_2$
under this exact white Gaussian model. Consequently the
periodogram variance does not vanish as $T$ increases.
Averaging approximately independent tapered or segmented
estimates reduces variance, at a cost of frequency resolution
and possible bias. For $K$ independent identical exponential
spectral ordinates, $2K\bar I/S\sim\chi^2_{2K}$;
Welch overlap and tapered spectra violate exact independence,
so their effective degrees of freedom need calibration.

\section{Short-time Fourier transform, wavelets and admissibility}
A window $w$ produces the short-time Fourier transform
$V_wx(\tau,f)=\int x(t)\overline{w(t-\tau)}
 e^{-2\pi ift}\,dt$. Window duration sets an uncertainty
trade-off; the display resolution is not a ground-truth
measurement of a neural oscillation's instantaneous frequency.
For a continuous wavelet $\psi$, the usual admissibility
condition is
\[
C_\psi=\int_{\mathbb R}\frac{|\Psi(f)|^2}{|f|}\,df<\infty.
\]
It requires zero mean for sufficiently regular wavelets,
so an uncorrected finite-frequency Morlet wavelet is only
approximately admissible; a DC correction makes the distinction
explicit. Reconstruction follows with the appropriate wavelet
normalization. Discrete wavelet frames require bounds
$A\|x\|^2\le\sum_j|\langle x,\psi_j\rangle|^2
\le B\|x\|^2$ for positive finite $A,B$.

\section{Coherence, phase consistency and finite samples}
Even when the population coherence is zero, the squared
magnitude of a finite-sample coherence estimate is nonnegative
and generally biased upward. A simple idealized average of
$K$ independent complex Gaussian spectral replicates gives,
under zero population coherence, an estimator with
$\E[\widehat C]\approx1/K$; this is \emph{not} a universal
threshold for overlapped, tapered or strongly dependent data.
For iid spike phases $\phi_1,\ldots,\phi_N$, the usual
phase-locking value (PLV) is $|N^{-1}\sum_j e^{i\phi_j}|$.
Even uniformly distributed phases produce a positive finite-sample
value. The pairwise phase consistency (PPC)
\[
\widehat{\mathrm{PPC}}=
\frac{1}{N(N-1)}\sum_{j\ne k}
\cos(\phi_j-\phi_k)
\]
has zero expectation for independent uniform phases,
provided the sampling and independence conditions hold.
Surrogates made by Fourier phase randomization typically preserve
individual Fourier magnitudes; iterative amplitude-adjusted Fourier
transforms (IAAFT) approximately preserve both the marginal amplitude
distribution and a target power spectrum. Neither automatically
preserves cross-channel dependence or experimental trial structure.

\begin{example}[What is—and is not—implemented]\sloppy
\path{welch_psd.m},
\path{multitaper_psd.m},
\path{wavelet_coherence.m} and
\path{pairwise_phase_consistency.m}
produce numerical estimators. A scientific significance statement
requires effective degrees of freedom, trial/block design and
an appropriate null; it is not guaranteed by the function name.
\end{example}
% --- expanded edition (October 2026) ---
\section{Morlet wavelets: resolution and normalization}
A complex Morlet wavelet centred at frequency $f_0$ with $n_c$ cycles is
\[
\psi(t)=A\exp(2\pi if_0t)\exp\bigl(-t^2/(2\sigma_t^2)\bigr),\qquad
\sigma_t=\frac{n_c}{2\pi f_0}.
\] Its Fourier transform is a Gaussian centred at
$f_0$ with standard deviation $\sigma_f=1/(2\pi\sigma_t)=f_0/n_c$, so
that $\sigma_t\sigma_f=1/(2\pi)$. The energy densities $|\psi|^2$ and
$|\Psi|^2$ have standard deviations $\sigma_t/\sqrt2$ and
$\sigma_f/\sqrt2$, whose product $1/(4\pi)$ is the minimum allowed by the
uncertainty principle (Result 61), as for every Gaussian. The \emph{number of cycles}
is the single parameter that sets the trade-off: at $f_0=40$\,Hz with
$n_c=7$, $\sigma_t\approx28$\,ms and $\sigma_f\approx5.7$\,Hz; at
$f_0=4$\,Hz the same $n_c$ gives $\sigma_t\approx280$\,ms and
$\sigma_f\approx0.57$\,Hz. Constant $n_c$ is constant relative
bandwidth, the defining property of a wavelet transform.

Two normalizations are common. With $A=(\sigma_t\sqrt{\pi})^{-1/2}$ the
wavelet has unit energy and white noise yields equal expected power at
all scales. With $A$ chosen so that $|\Psi(f_0)|=1$, a unit-amplitude
sinusoid at $f_0$ yields unit amplitude, which is convenient for
envelope interpretation. Power spectra computed with the two conventions
differ by a factor proportional to $f_0$, which changes the apparent
$1/f$ slope by one. The convention must be reported.

\section{Testing phase concentration}
For $N$ independent phases, the Rayleigh statistic $Z=N R^2$, with
$R$ the mean resultant length, is approximately exponential with mean
one under uniformity; a refined approximation of the $p$-value is
\[
p\approx\exp\Bigl(\sqrt{1+4N+4(N^2-NZ)}-(1+2N)\Bigr).
\]
The test is most powerful against unimodal (von Mises) alternatives and
has no power against bimodal locking with opposite modes, for which
the resultant vanishes. The relation $\widehat{\mathrm{PPC}}=(NR^2-1)/(N-1)$
connects the PPC of this chapter to the Rayleigh statistic, so the two
measures carry the same information about a single set of phases; PPC
is preferred for \emph{comparing} conditions because its expectation does
not depend on $N$.

\begin{proposition}
For independent phases with common resultant $\rho=|\E e^{i\phi}|$,
$\E[\widehat{\mathrm{PPC}}]=\rho^2$.
\end{proposition}
\begin{proof}
For $j\ne k$, independence gives
$\E e^{i(\phi_j-\phi_k)}=\E e^{i\phi_j}\overline{\E e^{i\phi_k}}=\rho^2$,
which is real; its real part is $\E\cos(\phi_j-\phi_k)$. Average over the
$N(N-1)$ ordered pairs.
\end{proof}
Spikes from the same burst, or many spikes from one trial, violate
independence. The PPC variants that only pair spikes from \emph{different
trials} remove this dependence at the cost of variance.

\section{Why phase-randomized surrogates preserve the spectrum}
Given a real series of length $T$ with DFT $X_k$, a phase-randomized
surrogate replaces $X_k$ by $|X_k|e^{i\theta_k}$ with independent uniform
$\theta_k$ for $0<k<T/2$, sets $\theta_{T-k}=-\theta_k$ to keep the
result real, and inverts. Its periodogram is identical to the original's
by construction, and its circular autocovariance is therefore identical
too (the periodogram is the DFT of the circular sample autocovariance).
Under the null hypothesis that the data are a linear Gaussian stationary
process, the surrogates are draws from (approximately) the same
distribution; nonlinear statistics such as PAC, bicoherence or
time-asymmetry then have a valid reference distribution. The surrogate
destroys every non-Gaussian feature, including spikes, sharp waves and
bursts, so a rejection means only ``not a stationary linear Gaussian
process'', which is almost always true of brain signals.

\section*{Exercises for this chapter}
\begin{exercise}
Derive $\sigma_f=f_0/n_c$ for the Morlet wavelet from the Fourier
transform of a Gaussian.
\end{exercise}
\begin{exercise}
Show that the mean resultant length of phases concentrated equally at
$0$ and $\pi$ is zero, and propose a statistic that detects such
bimodal locking.
\end{exercise}
\begin{exercise}
Prove that the periodogram is the DFT of the circular sample
autocovariance.
\end{exercise}
\begin{exercise}
Simulate phases from a von Mises distribution with $\kappa=0.5$ and
compare the average PLV and PPC for $N=10$ and $N=1000$.
\end{exercise}
